% ============================================================================
% Formal Equivalence Proof: Box-Constraint CBF-QP = Clipping
% For SafeContract paper (ICRA Workshop 2026)
% ============================================================================
%
% This can be included as an appendix or inlined in the experiments section.
% The main paper already states the result at line 199 of main_full_6page.tex;
% this file provides the full formal treatment.

\begin{proposition}[Box-Constraint QP Equivalence]
\label{prop:qp-clip}
Let $\mathcal{C} = \{\mathbf{a} \in \mathbb{R}^d : l_i \leq a_i \leq u_i,\; \forall\, i\}$ be a hyperrectangular constraint set. The safety-filter QP
\begin{equation}
  \mathbf{a}_{\mathrm{safe}} = \arg\min_{\mathbf{a}' \in \mathcal{C}} \|\mathbf{a}' - \mathbf{a}\|_2^2
\end{equation}
has the closed-form solution
\begin{equation}
  \mathbf{a}_{\mathrm{safe}} = \mathrm{clip}(\mathbf{a},\; \mathbf{l},\; \mathbf{u}).
\end{equation}
\end{proposition}

\begin{proof}
The objective decomposes as
\[
  \|\mathbf{a}' - \mathbf{a}\|_2^2 = \sum_{i=1}^{d} (a'_i - a_i)^2,
\]
and the constraint set $\mathcal{C}$ is the Cartesian product $[l_1, u_1] \times \cdots \times [l_d, u_d]$.
Both the objective and the feasible set are separable across dimensions, so the
$d$-dimensional QP decomposes into $d$ independent scalar optimization problems:
\begin{equation}
  a_{\mathrm{safe},i} = \arg\min_{a'_i \in [l_i, u_i]} (a'_i - a_i)^2, \quad i = 1, \ldots, d.
  \label{eq:scalar-qp}
\end{equation}

Each scalar problem~\eqref{eq:scalar-qp} minimizes a strictly convex quadratic
over a closed interval. The unconstrained minimizer is $a'_i = a_i$. Three cases arise:

\begin{enumerate}
  \item If $a_i \in [l_i, u_i]$: the unconstrained minimizer is feasible, so $a_{\mathrm{safe},i} = a_i$.
  \item If $a_i < l_i$: the objective is strictly increasing on $[l_i, u_i]$, so the minimum is at $a_{\mathrm{safe},i} = l_i$.
  \item If $a_i > u_i$: the objective is strictly decreasing on $[l_i, u_i]$, so the minimum is at $a_{\mathrm{safe},i} = u_i$.
\end{enumerate}

These three cases are exactly the definition of the clipping function:
\begin{equation}
  a_{\mathrm{safe},i} = \max\bigl(l_i,\; \min(a_i,\; u_i)\bigr) = \mathrm{clip}(a_i,\; l_i,\; u_i).
\end{equation}

Combining across all dimensions:
\begin{equation}
  \mathbf{a}_{\mathrm{safe}} = \mathrm{clip}(\mathbf{a},\; \mathbf{l},\; \mathbf{u}). \qedhere
\end{equation}
\end{proof}

\paragraph{Implication for SafeContract vs.\ AEGIS.}
AEGIS~\cite{aegis} solves a CBF-based QP at each timestep to project unsafe actions
onto a safe set defined by control barrier functions. When the safe set is a
hyperrectangle (box constraints on joint positions, velocities, or end-effector
coordinates), Proposition~\ref{prop:qp-clip} shows that the QP solution reduces
to element-wise clipping. SafeContract's \texttt{enforce()} operation performs
exactly this clipping. Therefore:

\begin{itemize}
  \item \textbf{For box constraints}, SafeContract and AEGIS produce \emph{identical}
        safe actions, but SafeContract does so in ${\sim}13\,\mu$s without a QP solver,
        VLM, or depth sensor, compared to AEGIS's ${\sim}360\,\mu$s with full
        perception stack.
  \item \textbf{For non-box constraints} (e.g., obstacle avoidance CBFs, workspace
        polytopes, nonlinear barrier functions), AEGIS handles constraint geometries
        that SafeContract cannot express. SafeContract is limited to axis-aligned
        hyperrectangular constraints.
  \item \textbf{The approaches are complementary}: SafeContract provides fast
        action-space enforcement plus violation analytics (fingerprinting, shift
        detection), while AEGIS provides scene-aware spatial safety. A deployment
        could layer both - SafeContract for microsecond action-space bounds, AEGIS
        for obstacle avoidance when scene perception is available.
\end{itemize}

\paragraph{Comparison paragraph (for Related Work / Discussion).}
% Copy this into the paper body:
%
% For box constraints, SafeContract's clipping produces outputs mathematically
% identical to AEGIS's CBF-QP solution, because the QP for hyperrectangular
% constraints has a closed-form clip solution (Proposition~\ref{prop:qp-clip}).
% The difference is scope: AEGIS handles arbitrary CBF constraints - obstacle
% avoidance, scene-aware safety - at the cost of a VLM, depth sensor, and QP
% solver (${\sim}360\,\mu$s). SafeContract handles action-space constraints only,
% requiring no external perception (${\sim}13\,\mu$s). Where AEGIS and Sentinel
% offer scene-level reasoning that SafeContract cannot, SafeContract provides
% violation fingerprinting and statistical shift detection that neither offers.
% The approaches are complementary: SafeContract enforces fast action-space bounds
% while AEGIS or Sentinel can layer scene-aware constraints on top.
